\exercise{9.8}{Finite and infinite possibilities for five limits}

\goal Determine each of the five requested limits, using the
convention introduced in Definition 9.8: an existing limit may be
a real number, $+\infty$, or $-\infty$. The answer ``NOT EXIST''
therefore means that none of these three possibilities works.

\idea Use the defining threshold tests. A sequence has limit
$+\infty$ if, for every $M>0$, all sufficiently late terms exceed
$M$. It has limit $-\infty$ if, for every $M<0$, all sufficiently
late terms are below $M$. Neither definition describes finite
convergence, and the finite-limit arithmetic laws cannot simply
be applied to expressions involving infinity.

\step{(a) The sequence $n^3$.}
Let $M>0$ be arbitrary. Set $N=\max\{1,M\}$. If $n>N$, then
$n\geq1$ and
\[
 n^3=n\cdot n^2\geq n>M.
\]
Thus every prescribed positive threshold is exceeded by all
sufficiently late terms, exactly as Definition 9.8 requires.

\answer{In part (a), $\displaystyle\lim_{n\to\infty}n^3=+\infty$.}

\step{(b) The sequence $-n^3$.}
Let $M<0$. The number $-M$ is positive. Choose
$N=\max\{1,-M\}$. If $n>N$, the estimate used in (a) gives
$n^3\geq n>-M$. Multiplication by $-1$ reverses the inequality,
giving $-n^3<M$. This is the negative-infinity definition.

\answer{In part (b), $\displaystyle\lim_{n\to\infty}(-n^3)=-\infty$.}

\step{(c) The sequence $(-n)^n$.}
At an even index $n=2m$ the value is $(2m)^{2m}>0$; at an
odd index $n=2m+1$ it is $-(2m+1)^{2m+1}<0$.
There are even and odd indices beyond every threshold: choose
a positive integer $m>N$ to make both $2m$ and $2m+1$ exceed
a prescribed $N$.

The limit cannot be $+\infty$, because the test with $M=1$
would require all sufficiently late terms to be greater than $1$,
whereas arbitrarily late odd terms are negative. It cannot be
$-\infty$, because the test with $M=-1$ would require all
sufficiently late terms to be less than $-1$, whereas arbitrarily
late even terms are positive.

It also cannot be a finite real number $L$. If it converged to
$L$, the definition with $\varepsilon=1$ would give a threshold
$N_1$ such that
\[
 n>N_1\quad\Longrightarrow\quad\abs{(-n)^n-L}<1.
\]
Choose an even positive integer $n>N_1$ with
$n>\max\{1,L+1\}$. Such an index exists by taking $n=2m$
with $m$ sufficiently large. Since $n\geq1$, $n^n\geq n$,
so at this even index
\[
 (-n)^n=n^n\geq n>L+1.
\]
Consequently $\abs{(-n)^n-L}=n^n-L>1$, contradicting the
required error bound. All three possible types of limit have
now been excluded.

\answer{In part (c), the limit does NOT EXIST, even in the
extended sense of Definition 9.8.}

\step{(d) The sequence $(1.01)^n$.}
Write $1.01=1+1/100$. The binomial theorem, or Bernoulli's
inequality, gives
\[
 \left(1+\frac1{100}\right)^n\geq1+\frac n{100}.
\]
For completeness, the inequality $(1+x)^n\geq1+nx$ for
$x\geq0$ follows by induction: the case $n=1$ is equality,
and from the inequality at $n$ one obtains
\[
 (1+x)^{n+1}\geq(1+nx)(1+x)
 =1+(n+1)x+nx^2\geq1+(n+1)x.
\]
Now let $M>0$ and set $N=100M$. If $n>N$, then
$n/100>M$, and hence
\[
 (1.01)^n\geq1+\frac n{100}>M.
\]
This proves divergence to positive infinity by the definition.

\answer{In part (d), $\displaystyle\lim_{n\to\infty}(1.01)^n=+\infty$.}

\step{(e) The sequence $n^n$.}
For $n\geq1$, the product of $n$ factors all equal to $n$
satisfies $n^n\geq n$. Let $M>0$ and set
$N=\max\{1,M\}$. If $n>N$, then
$n^n\geq n>M$. The definition therefore applies.

\answer{In part (e), $\displaystyle\lim_{n\to\infty}n^n=+\infty$.}

\exercise{9.9}{Order preserved by existing extended limits}

\goal Assume there is a threshold $N_0$ such that
$s_n\leq t_n$ whenever $n>N_0$. Prove the three assertions,
including every possible finite or infinite case in part (c).

\idea Combine the threshold at which the comparison becomes
valid with the threshold supplied by the relevant limit. In the
finite case, two limits in the wrong order would force the late
terms to have the wrong order as well.

\step{(a) From $s_n\to+\infty$ to $t_n\to+\infty$.}
Fix $M>0$. Since $s_n\to+\infty$, some $N_1$ satisfies
$s_n>M$ for all $n>N_1$. Set $N=\max\{N_0,N_1\}$.
For every $n>N$, both the comparison and this threshold test
hold, so
\[
 t_n\geq s_n>M.
\]
This holds for the arbitrary $M>0$, proving $t_n\to+\infty$.

\step{(b) From $t_n\to-\infty$ to $s_n\to-\infty$.}
Fix $M<0$. Since $t_n\to-\infty$, some $N_1$ satisfies
$t_n<M$ for all $n>N_1$. With
$N=\max\{N_0,N_1\}$, every $n>N$ gives
\[
 s_n\leq t_n<M.
\]
This is the definition of $s_n\to-\infty$.

\step{(c) Account for the infinite limit cases.}
Write $S=\lim s_n$ and $T=\lim t_n$, which exist by the
assumption in this part. Their possible values lie in
$\R\cup\{-\infty,+\infty\}$, with the usual extended order.
The finite and infinite alternatives cannot overlap. Indeed,
convergence to a real $c$ puts the late terms between $c-1$
and $c+1$. A simultaneous limit $+\infty$ would put them
above $\max\{1,c+2\}$, contradicting the upper restriction;
a simultaneous limit $-\infty$ would put them below
$\min\{-1,c-2\}$, contradicting the lower restriction.
Limits of both signs of infinity are incompatible as well,
since they would eventually require a term to be both greater
than $1$ and less than $-1$.
The following cases exhaust all possibilities involving infinity:
\begin{enumerate}
 \item If $S=-\infty$, then $S\leq T$ for every possible $T$.
 \item If $T=+\infty$, then $S\leq T$ for every possible $S$.
 \item If $S=+\infty$, part (a) forces $T=+\infty$, so
       $S=T$ and in particular $S\leq T$.
 \item If $T=-\infty$, part (b) forces $S=-\infty$, so
       again $S=T$.
\end{enumerate}
After these cases, both $S$ and $T$ must be finite real numbers.

\step{(c) Prove the finite case by contradiction.}
Suppose, contrary to the desired inequality, that $S>T$.
Set $\delta=(S-T)/3>0$. By finite convergence there are
thresholds $N_1,N_2$ such that
\[
 n>N_1\Longrightarrow\abs{s_n-S}<\delta,
 \qquad
 n>N_2\Longrightarrow\abs{t_n-T}<\delta.
\]
Thus $s_n>S-\delta$ and $t_n<T+\delta$ at the corresponding
late indices. But
\[
 (S-\delta)-(T+\delta)
 =(S-T)-2\delta
 =(S-T)-\frac{2(S-T)}3
 =\frac{S-T}{3}>0.
\]
For $n>\max\{N_0,N_1,N_2\}$ this gives
\[
 s_n>S-\delta>T+\delta>t_n,
\]
contradicting $s_n\leq t_n$. Hence $S\leq T$ in the finite
case as well.

\answer{The eventual order implies both infinite-limit comparisons
in (a) and (b), and $\lim s_n\leq\lim t_n$ whenever both limits
exist in the extended sense.}

\exercise{9.10}{Scaling and reversing an infinite limit}

\goal Prove the effects of multiplying a sequence with limit
$+\infty$ by a positive constant or a negative constant, and
prove both directions of the sign-reversal equivalence.

\idea Choose a threshold for $s_n$ by dividing the target
threshold by the multiplying constant. A negative constant
reverses the inequality, so the signs must be recorded.

\step{(a) Multiplication by $k>0$.}
Assume $s_n\to+\infty$ and fix $k>0$.
To prove $ks_n\to+\infty$, let $M>0$. Then $M/k>0$,
so the assumed limit supplies $N$ with
\[
 n>N\quad\Longrightarrow\quad s_n>\frac Mk.
\]
Multiplying by $k>0$ preserves the direction and gives
$ks_n>M$. This verifies the target definition.

In fact this implication is reversible. If $ks_n\to+\infty$
and $M>0$, use the positive threshold $kM$ for $ks_n$.
Eventually $ks_n>kM$, and division by $k>0$ gives $s_n>M$.
Thus positive multiplication preserves the $+\infty$ limit in
both directions.

\step{(b) Prove $s_n\to+\infty\Longrightarrow -s_n\to-\infty$.}
Let $M<0$ be an arbitrary threshold for $-s_n$.
Because $-M>0$, the assumption $s_n\to+\infty$ supplies
$N$ such that $n>N$ implies $s_n>-M$.
Multiplication by $-1$ reverses the inequality, giving
$-s_n<M$. Thus $-s_n\to-\infty$.

\step{(b) Prove the reverse implication.}
Assume $-s_n\to-\infty$ and let $M>0$.
The number $-M$ is negative, so some $N$ satisfies
$n>N\Longrightarrow -s_n<-M$. Multiplication by $-1$
reverses this inequality and gives $s_n>M$.
Consequently $s_n\to+\infty$, completing the equivalence.

\step{(c) Multiplication by $k<0$.}
Assume $s_n\to+\infty$ and $k<0$. Let $M<0$ be
an arbitrary target threshold for $ks_n$. The quotient
$M/k$ is positive, since both numerator and denominator are
negative. Hence the assumed limit supplies $N$ such that
$n>N$ implies $s_n>M/k$. Multiplying by $k<0$ reverses
the comparison:
\[
 n>N\quad\Longrightarrow\quad ks_n<M.
\]
This proves $ks_n\to-\infty$.

The corresponding reverse implication also holds. If
$ks_n\to-\infty$, fix $M>0$ and use the negative threshold
$kM$. Eventually $ks_n<kM$, and division by $k<0$ reverses
the inequality, yielding $s_n>M$. Therefore $s_n\to+\infty$.

\answer{Positive scaling preserves a positive-infinity limit;
negative scaling changes it to negative infinity. In particular,
$s_n\to+\infty$ if and only if $-s_n\to-\infty$, with each
direction justified by its own threshold test.}

\exercise{9.11}{Adding a sequence that is bounded below}

\goal Prove the sum has limit $+\infty$ in each of the three
stated situations. Distinguish a lower bound valid for all terms
from a lower bound valid only after an initial segment.

\idea Once $t_n\geq m$ holds in the tail, it suffices to make
$s_n>M-m$. If this number is not positive, replace it by a
larger positive threshold before invoking Definition 9.8.

\step{(a) Use the stated global lower bound.}
Assume $s_n\to+\infty$ and
$\inf\{t_n:n\in\N\}>-\infty$.
The set is nonempty and consists of real numbers; its infimum
is at most the finite member $t_1$. Thus this hypothesis provides
a real lower bound $m$, for example the real infimum itself,
such that $t_n\geq m$ for every $n$.

Let $M>0$ be arbitrary, and define
\[
 B=\max\{1,M-m\}>0.
\]
Since $s_n\to+\infty$, there is $N$ such that $n>N$
implies $s_n>B$. For those indices,
\[
 s_n+t_n\geq s_n+m>B+m\geq(M-m)+m=M.
\]
The first inequality uses the global lower bound; the last uses
$B\geq M-m$. This proves the positive-infinity limit.

\step{(b) Derive an eventual lower bound in every allowed limit case.}
Assume $s_n\to+\infty$ and $T=\lim t_n>-\infty$.
The existing limit $T$ is either a real number or $+\infty$.
\begin{itemize}
 \item If $T\in\R$, apply finite convergence with
       $\varepsilon=1$. For some $N_1$, every $n>N_1$
       has $\abs{t_n-T}<1$, implying $t_n>T-1$.
       Thus we may take $m=T-1$ as an eventual lower bound.
 \item If $T=+\infty$, use the positive threshold $1$ in
       Definition 9.8. For some $N_1$, every $n>N_1$
       has $t_n>1$. Thus we may take $m=1$.
\end{itemize}
In either case there are a real $m$ and a threshold $N_1$ such
that $t_n\geq m$ whenever $n>N_1$. This is a tail bound;
no assertion about the earlier terms is needed.

Now fix $M>0$. Set $B=\max\{1,M-m\}>0$, and use
$s_n\to+\infty$ to choose $N_2$ such that $s_n>B$
for $n>N_2$. For
$n>\max\{N_1,N_2\}$, both estimates hold and
\[
 s_n+t_n\geq s_n+m>B+m\geq M.
\]
This proves the desired limit using the eventual bound directly,
without applying a global hypothesis to an unexamined prefix.

\step{(c) A bounded sequence is bounded below.}
If $(t_n)$ is bounded, there is a real $C\geq0$ with
$\abs{t_n}\leq C$ for every $n$. Therefore $t_n\geq-C$
for every $n$, a global lower bound, and part (a) applies.
One can also see the threshold explicitly: for $M>0$, choose
$N$ such that $s_n>M+C$ whenever $n>N$; this is a valid
positive threshold because $M+C>0$. Then
\[
 s_n+t_n\geq s_n-C>M+C-C=M.
\]

\answer{In (a), (b), and (c), the sum $s_n+t_n$ tends to
$+\infty$. An eventual lower bound on the added sequence already
suffices, as the proof of (b) shows.}

\exercise{9.12}{The ratio test for the size of a sequence}

\goal Suppose $s_n\ne0$ for every $n$, and the limit
\[
 L=\lim_{n\to\infty}r_n,
 \qquad r_n=\abs{\frac{s_{n+1}}{s_n}},
\]
exists. Prove $s_n\to0$ if $L<1$, and
$\abs{s_n}\to+\infty$ if $L>1$, including the case
$L=+\infty$.

\idea A ratio eventually less than a fixed $q<1$ gives a
geometric bound. For ratios tending to a value larger than $1$,
apply that first result to the reciprocal magnitudes.

\step{First locate the possible values of $L$.}
Every $r_n$ is positive, since its numerator and denominator are
nonzero in absolute value. Therefore its existing limit can only
be a nonnegative real number or $+\infty$.
For example, if a finite limit were $L<0$, convergence with
$\varepsilon=-L/2>0$ would eventually give
$r_n<L+(-L/2)=L/2<0$, a contradiction. A limit $-\infty$
would eventually give $r_n<-1$ by its definition, also a
contradiction. This verifies the claimed restriction on $L$.

\step{(a) Choose a single ratio bound valid at every sufficiently late index.}
Suppose $L<1$. Then $0\leq L<1$, so
\[
 q=\frac{L+1}{2}
 \quad\hbox{satisfies}\quad 0<q<1\quad\hbox{and}\quad L<q.
\]
By convergence of $r_n$ to the finite number $L$, some threshold
$N_0$ gives $\abs{r_n-L}<q-L$ for all $n>N_0$.
In particular $r_n<q$ at all those indices.
Choose an integer $N\in\N$ with $N>N_0$.
Then every integer $n\geq N$ has $n>N_0$, so
\[
 \abs{s_{n+1}}=r_n\abs{s_n}<q\abs{s_n}
 \qquad(n\geq N).
\]
The integer choice makes the starting index and the validity of
every subsequent ratio inequality explicit.

\step{(a) Prove the geometric estimate by induction.}
For $n=N+1$, the preceding inequality gives
\[
 \abs{s_{N+1}}<q\abs{s_N}=q^{(N+1)-N}\abs{s_N}.
\]
Assume for some integer $n>N$ that
$\abs{s_n}<q^{n-N}\abs{s_N}$. Since the ratio estimate
also applies at this $n$, multiplication by $q>0$ gives
\[
 \abs{s_{n+1}}<q\abs{s_n}
 <q\,q^{n-N}\abs{s_N}
 =q^{(n+1)-N}\abs{s_N}.
\]
Induction proves
\[
 \abs{s_n}<q^{n-N}\abs{s_N}\qquad(n>N).
\]

Since $0<q<1$, Theorem 9.7(b) gives $q^m\to0$ as
$m\to\infty$. To spell out the use of this fact, let
$\varepsilon>0$. The number $C=\abs{s_N}$ is positive,
because $s_N\ne0$. Choose a threshold $K$ such that
$m>K$ implies $q^m<\varepsilon/C$.
If $n>N+\max\{K,0\}$, then $m=n-N>K$ and $n>N$,
so
\[
 \abs{s_n-0}<q^{n-N}C<\frac{\varepsilon}{C}C
 =\varepsilon.
\]
This is precisely finite convergence to $0$.

\step{(b) Form the reciprocal sequence and compute its ratio.}
Suppose $L>1$, and define
\[
 t_n=\frac1{\abs{s_n}}>0.
\]
It is defined and nonzero for every $n$. Its absolute ratio is
\begin{align*}
 \abs{\frac{t_{n+1}}{t_n}}
 &=\frac{1/\abs{s_{n+1}}}{1/\abs{s_n}}\\
 &=\frac{\abs{s_n}}{\abs{s_{n+1}}}
 =\frac1{r_n}.
\end{align*}
There are two distinct cases for the ratio limit:
\begin{itemize}
 \item If $1<L<+\infty$, the finite reciprocal law applies
       to $r_n\to L$, since every $r_n\ne0$ and $L\ne0$.
       It gives $1/r_n\to1/L$, with $0<1/L<1$.
 \item If $L=+\infty$, every $r_n$ is positive, so
       Theorem 9.10 gives $1/r_n\to0$. This conclusion uses
       the infinite-limit reciprocal theorem, not the finite
       quotient law with an infinite denominator.
\end{itemize}
In either case, the absolute ratio for $t_n$ has a finite limit
strictly less than $1$. Part (a) therefore proves $t_n\to0$.
Finally apply Theorem 9.10 to the positive sequence
$\abs{s_n}$: its reciprocal is $t_n\to0$, so
\[
 \abs{s_n}\longrightarrow+\infty.
\]

\answer{If $L<1$, the sequence itself converges to $0$.
If $L>1$, including $L=+\infty$, its magnitudes tend to
$+\infty$. No conclusion for the boundary value $L=1$ is
asserted by either part.}

\exercise{9.13}{All real bases in a geometric sequence}

\goal Prove the complete classification of $a^n$ as the fixed
real base $a$ varies, including $a=0$, $a=1$, $a=-1$, and
both intervals of negative bases.

\idea Use Exercise 9.12 where the terms are nonzero. For a
negative base with magnitude at least $1$, check that no finite
limit and neither infinite limit can accommodate the alternating
signs.

\step{Case 1: $a=0$, and then $0<\abs a<1$.}
If $a=0$, every term $a^n$ is $0$ because $n\geq1$.
Thus the sequence converges to $0$ directly.
If $0<\abs a<1$, every term is nonzero and
\[
 \abs{\frac{a^{n+1}}{a^n}}=\abs a<1.
\]
This constant ratio has limit $\abs a$, so Exercise 9.12(a)
proves $a^n\to0$. The two cases together cover all
$\abs a<1$, including negative $a$ of magnitude less than $1$.

\step{Case 2: $a=1$.}
Every term equals $1^n=1$. Thus for every $\varepsilon>0$,
$\abs{a^n-1}=0<\varepsilon$ at every index, giving limit $1$.

\step{Case 3: $a>1$.}
The terms $a^n$ are positive and nonzero. Their absolute ratio
is the constant $a>1$. Exercise 9.12(b) gives
$\abs{a^n}\to+\infty$. Since $a^n>0$,
$\abs{a^n}=a^n$, so the sequence itself tends to $+\infty$.

\step{Case 4: $a=-1$.}
The terms are $1$ at even indices and $-1$ at odd indices.
If there were a finite limit $L$, convergence with
$\varepsilon=1/2$ would put every sufficiently late term
within $1/2$ of $L$. Choose a late even index $p$ and a late
odd index $q$. The triangle inequality would give
\[
 2=\abs{1-(-1)}
 \leq\abs{a^p-L}+\abs{L-a^q}<\frac12+\frac12=1,
\]
a contradiction. A limit $+\infty$ fails the threshold $M=1$
because late odd terms equal $-1$. A limit $-\infty$ fails
the threshold $M=-1$ because late even terms equal $1$.
Thus no extended limit exists.

\step{Case 5: $a<-1$.}
Set $b=\abs a=-a>1$. By Case 3, $b^n\to+\infty$.
The terms satisfy
\[
 a^n=(-1)^n b^n,\qquad\abs{a^n}=b^n.
\]
Suppose first that $a^n$ had a finite limit $L$.
Convergence with $\varepsilon=1$ would give, eventually,
\[
 \abs{a^n}\leq\abs{a^n-L}+\abs L<1+\abs L.
\]
But $b^n\to+\infty$ means that eventually
$b^n>1+\abs L$, contradicting $\abs{a^n}=b^n$.
Thus there is no finite limit.

There is also no limit $+\infty$, since for every threshold
index there is a later odd index with $a^n<0$, defeating
the positive threshold $M=1$. There is no limit $-\infty$,
since there is always a later even index with $a^n>0$,
defeating the negative threshold $M=-1$.
The required even and odd indices can be taken as $2m$ and
$2m+1$, where $m$ is any positive integer larger than the
proposed threshold index.

\answer{For every real $a$,
\[
 \lim_{n\to\infty}a^n=
 \begin{cases}
  0,&\abs a<1,\\
  1,&a=1,\\
  +\infty,&a>1,\\
  \text{does not exist},&a\leq-1.
 \end{cases}
\]
The last line excludes finite limits and both infinite limits.}

\exercise{9.14}{Geometric growth compared with a real power}

\goal Fix $p>0$ and determine the limit of $a^n/n^p$
for every real $a$, using Exercise 9.12 where its hypotheses
apply. In particular, include both boundary bases $a=1$ and
$a=-1$.

\idea The absolute ratio is
$\abs a\,(n/(n+1))^p$. Prove its second factor tends to $1$,
then distinguish ratios below $1$, above $1$, and the boundary
cases where the ratio test does not decide the answer.

\step{1. State the real-power facts and justify the needed limit.}
As in Theorem 9.7(a), positive bases with real exponents use
the usual real-power algebra and order: products and quotients
obey the exponent rules; and $r^v\leq r^u$ whenever
$0<r<1$ and $0<u<v$. These are the elementary
real-power properties already implicit in the statement of
Theorem 9.7(a). We do not assume a continuity theorem for
real powers from a later section.

Set $r_n=n/(n+1)$, so $0<r_n<1$. Choose an integer $m>p$,
possible by the Archimedean property. Then
\[
 r_n^m\leq r_n^p\leq1,
 \qquad
 0\leq1-r_n^p\leq1-r_n^m.
\]
For the integer exponent $m$, finite multiplication gives the
factorization
\[
 1-r_n^m=(1-r_n)(1+r_n+\cdots+r_n^{m-1}).
\]
Every one of the $m$ terms in the parenthesized sum is at
most $1$, and $1-r_n=1/(n+1)$. Therefore
\[
 0\leq1-r_n^p\leq1-r_n^m
 \leq\frac{m}{n+1}.
\]
The final expression tends to $0$: for $\varepsilon>0$,
$n>m/\varepsilon$ makes $m/(n+1)<\varepsilon$.
This direct error estimate proves
\[
 \left(\frac n{n+1}\right)^p\longrightarrow1
\]
for the possibly nonintegral exponent $p>0$.

\step{2. Compute the ratio when $a\ne0$.}
Let $u_n=a^n/n^p$. Since $n>0$, its denominator is positive.
If $a\ne0$, all $u_n$ are nonzero, and the real-power
quotient rule yields
\begin{align*}
 \abs{\frac{u_{n+1}}{u_n}}
 &=\abs{\frac{a^{n+1}}{(n+1)^p}\cdot\frac{n^p}{a^n}}\\
 &=\abs a\,\frac{n^p}{(n+1)^p}\\
 &=\abs a\left(\frac n{n+1}\right)^p.
\end{align*}
Step 1 and the finite scalar-multiple limit theorem show that
this ratio has limit $\abs a$.

\step{3. Treat $a=0$ and the cases $0<\abs a<1$.}
For $a=0$, every numerator is $0$ and every denominator is
positive, so $u_n=0$ and the limit is $0$.
For $0<\abs a<1$, Step 2 gives an absolute ratio limit
less than $1$, so Exercise 9.12(a) applies and gives $u_n\to0$.

\step{4. Treat the two boundary cases $\abs a=1$.}
If $a=1$ or $a=-1$, then $\abs{a^n}=1$, and hence
\[
 \abs{u_n-0}=\frac1{n^p}.
\]
Theorem 9.7(a), whose hypothesis is precisely $p>0$, says
$1/n^p\to0$. Thus $u_n\to0$ in both boundary cases.
The ratio limit here is $1$; the conclusion comes from the
direct error bound rather than from an inapplicable strict
case of Exercise 9.12.

\step{5. Treat $a>1$.}
The absolute ratio limit is $a>1$, so Exercise 9.12(b) gives
$\abs{u_n}\to+\infty$. Since $a^n>0$ and $n^p>0$,
the terms themselves are positive: $u_n=\abs{u_n}$.
Consequently $u_n\to+\infty$.

\step{6. Treat $a<-1$ and exclude all three types of limit.}
Exercise 9.12(b) gives $\abs{u_n}\to+\infty$, since
the absolute ratio limit is $\abs a>1$.
If $u_n$ converged to a finite $L$, then eventually
\[
 \abs{u_n}\leq\abs{u_n-L}+\abs L<1+\abs L
\]
by the error bound $\abs{u_n-L}<1$. This contradicts
$\abs{u_n}\to+\infty$.
At even indices $u_n>0$, while at odd indices $u_n<0$.
Every tail has an odd index, so the positive-infinity test
with $M=1$ fails. Every tail also has an even index, so
the negative-infinity test with $M=-1$ fails.
Thus no extended limit exists in this last case.

\answer{For $p>0$ and real $a$,
\[
 \lim_{n\to\infty}\frac{a^n}{n^p}=
 \begin{cases}
  0,&\abs a\leq1,\\
  +\infty,&a>1,\\
  \text{does not exist},&a<-1.
 \end{cases}
\]
In particular, the base $a=-1$ now gives limit $0$, unlike
the unscaled sequence in Exercise 9.13.}

\exercise{9.15}{Every fixed geometric sequence is dominated by a factorial}

\goal Prove $a^n/n!\to0$ for every $a\in\R$, including
zero and negative values.

\idea After the zero-base case is separated, the absolute
ratio is $\abs a/(n+1)$, which tends to $0$. This gives
the strict ratio condition from Exercise 9.12(a).

\step{1. Handle $a=0$ without dividing by a zero term.}
When $a=0$ and $n\geq1$, $a^n=0$. Since $n!>0$,
every term is $0$, and the limit is $0$.

\step{2. Compute the absolute ratio for $a\ne0$.}
Fix a nonzero real $a$ and write $s_n=a^n/n!$.
Then every $s_n$ is nonzero. By the factorial identity
$(n+1)!=(n+1)n!$,
\begin{align*}
 \abs{\frac{s_{n+1}}{s_n}}
 &=\abs{\frac{a^{n+1}}{(n+1)!}\cdot\frac{n!}{a^n}}\\
 &=\abs{\frac{a\,n!}{(n+1)n!}}\\
 &=\frac{\abs a}{n+1}.
\end{align*}
For $\varepsilon>0$, choosing $N>\abs a/\varepsilon$
makes $n>N$ imply $\abs a/(n+1)<\varepsilon$.
Thus the ratio limit is $L=0<1$.
Exercise 9.12(a) now proves $s_n\to0$. Its hypothesis is
about absolute ratios, so this argument applies to negative
$a$ as well as positive $a$.

\answer{For every real $a$, $\displaystyle\lim_{n\to\infty}a^n/n!=0$.}

\exercise{9.16}{Three applications of the infinite-limit results}

\goal Establish positive-infinity limits for all three given
expressions, using the proved infinite-limit theorems or
Exercises 9.9--9.15 with their hypotheses checked.

\idea In (a), isolate a factor with limit $+\infty$ and
another with a positive finite limit. In (b) and (c), use
Exercise 9.14 for the growing term and then control the
added term from below.

\step{(a) The expression $(n^4+8n)/(n^2+9)$.}
Every denominator $n^2+9$ is positive. Factor the expression
exactly as
\[
 \frac{n^4+8n}{n^2+9}
 =n^2\,\frac{1+8/n^3}{1+9/n^2}.
\]
Indeed, multiplying the numerator and denominator of the
right-hand fraction by $n^2$ gives
\[
 n^2\,\frac{n^2+8/n}{n^2+9}
 =\frac{n^4+8n}{n^2+9}.
\]
Theorem 9.7(a) gives $1/n^3\to0$ and $1/n^2\to0$.
The finite scalar and sum laws therefore give
$1+8/n^3\to1$ and $1+9/n^2\to1$.
All the latter denominators are positive, and the limit
denominator is $1\ne0$, so the finite quotient theorem
applies to this fraction and gives
\[
 v_n=\frac{1+8/n^3}{1+9/n^2}\longrightarrow1>0.
\]
Also $n^2\to+\infty$: for each $M>0$, every
$n>\sqrt M$ has $n^2>M$. Thus Theorem 9.9 applies to
$s_n=n^2$ and $t_n=v_n$, since its first limit is
$+\infty$ and its second limit is the positive number $1$.
It gives the desired product limit.

\answer{In part (a),
$\displaystyle\lim_{n\to\infty}\frac{n^4+8n}{n^2+9}=+\infty$.}

\step{(b) The expression $2^n/n^2+(-1)^n$.}
Exercise 9.14 with $a=2>1$ and $p=2>0$ gives
\[
 u_n=\frac{2^n}{n^2}\longrightarrow+\infty.
\]
The second sequence is bounded: $-1\leq(-1)^n\leq1$
for every $n$. Exercise 9.11(c) therefore applies to the
sum. In explicit threshold form, given $M>0$, choose $N$
so that $n>N$ implies $u_n>M+1$. Then
\[
 \frac{2^n}{n^2}+(-1)^n\geq u_n-1>M.
\]
This verifies the limit without requiring the oscillating
sequence $(-1)^n$ itself to have a limit.

\answer{In part (b),
$\displaystyle\lim_{n\to\infty}\left[\frac{2^n}{n^2}+(-1)^n\right]=+\infty$.}

\step{(c) The expression $3^n/n^3-3^n/n!$.}
Exercise 9.14 with $a=3$, $p=3$ and Exercise 9.15 with
$a=3$ give, respectively,
\[
 u_n=\frac{3^n}{n^3}\longrightarrow+\infty,
 \qquad
 v_n=\frac{3^n}{n!}\longrightarrow0.
\]
Thus $-v_n\to0$ by the finite scalar-multiple theorem.
Exercise 9.11(b) applies to $u_n+(-v_n)$, since the second
limit $0$ is greater than $-\infty$.
To verify the threshold directly, fix $M>0$.
Choose $N_1$ such that $n>N_1$ implies $u_n>M+1$.
Choose $N_2$ such that $n>N_2$ implies
$\abs{v_n}<1$, using its convergence to $0$.
For $n>\max\{N_1,N_2\}$ we have $v_n<1$ and hence
\[
 u_n-v_n>u_n-1>M.
\]
This is a growing term minus a term tending to $0$;
the proof does not subtract two infinite limits.

\answer{In part (c),
$\displaystyle\lim_{n\to\infty}\left[\frac{3^n}{n^3}-\frac{3^n}{n!}\right]=+\infty$.}

\exercise{9.17}{A definition-only proof for the sequence of squares}

\goal Use only Definition 9.8 to prove $n^2\to+\infty$.

\idea Given the positive height $M$, solve $n^2>M$ by
requiring $n>\sqrt M$. Both sides are positive, so squaring
this comparison preserves its direction.

\step{Formal proof with the quantifiers explicit.}
Let $M>0$ be arbitrary, and set $N=\sqrt M$.
If $n\in\N$ and $n>N$, then $n>\sqrt M>0$.
Multiplying positive inequalities, or using the strictly
increasing square function on the positive numbers, gives
\[
 n^2>(\sqrt M)^2=M.
\]
More explicitly, $n>\sqrt M$ gives
$n^2>n\sqrt M>(\sqrt M)^2$ after multiplication first
by $n>0$ and then by $\sqrt M>0$.
We have therefore proved that for every $M>0$ there is
$N=\sqrt M$ such that every positive integer $n>N$
satisfies $n^2>M$. This is exactly Definition 9.8.

\answer{$\displaystyle\lim_{n\to\infty}n^2=+\infty$, with the
explicit threshold $N=\sqrt M$ for each $M>0$.}

\exercise{9.18}{Finite geometric sums and their limits}

\goal Prove the finite-sum identity in (a), then determine
the limits in (b), (c), and (d), including the base $a=1$
where the identity's quotient is not defined.

\idea Subtract $a$ times the finite sum from the original
sum. Every interior power cancels. After the identity is
proved, use the established geometric-sequence limit for
$\abs a<1$, and a direct lower comparison for $a\geq1$.

\step{(a) Prove the identity for every allowed base and every index.}
Fix a real $a\ne1$ and a positive integer $n$, and define
\[
 S_n=1+a+a^2+\cdots+a^n=\sum_{k=0}^{n}a^k.
\]
Here $a^0$ denotes the constant empty product $1$, including
when $a=0$. Multiplication and reindexing a finite sum give
\[
 aS_n=\sum_{k=0}^{n}a^{k+1}
 =\sum_{j=1}^{n+1}a^j
 =a+a^2+\cdots+a^n+a^{n+1}.
\]
The reindexing sets $j=k+1$, so the lower bound $k=0$
becomes $j=1$ and the upper bound $k=n$ becomes $j=n+1$.
Subtracting the two finite sums cancels exactly the terms
$a,a^2,\ldots,a^n$:
\begin{align*}
 S_n-aS_n
 &=(1+a+\cdots+a^n)-(a+\cdots+a^n+a^{n+1})\\
 &=1-a^{n+1}.
\end{align*}
The left side factors as $(1-a)S_n$. Since $a\ne1$,
$1-a\ne0$, and division by this nonzero number yields
\[
 S_n=\frac{1-a^{n+1}}{1-a}.
\]
The base and index were arbitrary, so this proves the identity
for every real $a\ne1$ and every positive integer $n$.
The cancellation also works for $n=0$, if that index is
included, since then the identity is $1=(1-a)/(1-a)$.

\step{(b) Pass to the limit for $\abs a<1$.}
Now fix $a$ with $\abs a<1$. Then $a\ne1$, so (a)
applies. Exercise 9.13 gives $a^n\to0$.
Since $a^{n+1}=a\,a^n$, multiplication by the fixed real
number $a$ gives $a^{n+1}\to a\cdot0=0$ by the finite
scalar-multiple theorem. The denominator $1-a$ is a fixed
nonzero real number. Consequently the finite sum and quotient
laws applied to the identity give
\[
 \lim_{n\to\infty}S_n=\frac{1-0}{1-a}=\frac1{1-a}.
\]
One can read the same conclusion directly as an error
estimate, avoiding any ambiguity about division:
\begin{align*}
 \abs{S_n-\frac1{1-a}}
 &=\abs{\frac{1-a^{n+1}-1}{1-a}}\\
 &=\frac{\abs a^{\,n+1}}{\abs{1-a}}\longrightarrow0.
\end{align*}
The denominator in this expression is positive and constant;
the numerator tends to $0$. This includes $a=0$, when every
sum is exactly $1$.

\answer{For part (b), if $\abs a<1$, then
$\displaystyle\lim_{n\to\infty}(1+a+\cdots+a^n)=1/(1-a)$.}

\step{(c) Substitute $a=1/3$ and simplify the fraction.}
The displayed sequence of sums is
\[
 1+\frac13+\frac19+\frac1{27}+\cdots+\frac1{3^n}
 =\sum_{k=0}^{n}\left(\frac13\right)^k.
\]
Since $0<1/3<1$, part (b) applies. Its limit is
\[
 \frac1{1-1/3}
 =\frac1{(3-1)/3}
 =\frac1{2/3}
 =\frac32.
\]

\answer{The limit in part (c) is $3/2$.}

\step{(d) Handle $a=1$ and $a>1$ explicitly.}
If $a=1$, the sum contains $n+1$ terms, each equal to $1$,
so $S_n=n+1$. For an arbitrary $M>0$, $n>M$ implies
$S_n=n+1>M$. Therefore $S_n\to+\infty$ in this case.
The quotient in (a) was not used at $a=1$, where its
denominator would be zero.

If $a>1$, every term $a^k$ for $k=0,1,\ldots,n$ is at
least $1$: the exponent $0$ gives $1$, and a positive
exponent gives a product of factors greater than $1$.
Adding the $n+1$ inequalities gives
\[
 S_n=\sum_{k=0}^{n}a^k\geq\sum_{k=0}^{n}1=n+1.
\]
Since $n+1\to+\infty$, Exercise 9.9(a) proves
$S_n\to+\infty$. Equivalently, the same threshold $N=M$
works directly: $n>N$ implies $S_n\geq n+1>M$.

\answer{For part (d), every base $a\geq1$ gives
$\displaystyle\lim_{n\to\infty}(1+a+\cdots+a^n)=+\infty$.}
